
\documentclass[12pt,a4paper]{article}

\usepackage[a4paper,margin=1in]{geometry}
\usepackage{amsmath}
\usepackage{amssymb}
\usepackage{graphicx}
\usepackage{xcolor}
\usepackage{hyperref}
\usepackage{listings}
\usepackage{longtable}
\usepackage{array}
\usepackage{enumitem}
\usepackage{titlesec}
\usepackage{fancyhdr}
\usepackage{booktabs}
\usepackage[T1]{fontenc}

\hypersetup{
    colorlinks=true,
    linkcolor=blue,
    urlcolor=blue
}

\pagestyle{fancy}
\fancyhf{}
\rhead{CSI Intrusion Detection Project}
\lhead{Environment and Kernel Setup}
\cfoot{\thepage}

\title{\textbf{CSI Intrusion Detection Project\\
Complete Environment, Package, and Jupyter Kernel Setup Guide}}
\author{Rajshekhar Bhatta}
\date{\today}

\lstset{
    basicstyle=\ttfamily\small,
    backgroundcolor=\color{gray!10},
    frame=single,
    breaklines=true,
    columns=fullflexible
}

\begin{document}

\maketitle
\tableofcontents
\newpage

\section{Project Overview}

This document explains the complete environment setup procedure for the ESP32 CSI Intrusion Detection project.

The project uses:

\begin{itemize}
    \item ESP32 CSI passive sensing
    \item Python-based CSI parsing
    \item CNN+LSTM deep learning model
    \item TensorFlow GPU acceleration
    \item Jupyter Notebook development workflow
    \item Real-time CSI inference dashboard
\end{itemize}

\section{Project Architecture}

\begin{verbatim}
ESP32-E Passive Sniffer
        ↓
Serial CSI Stream
        ↓
Python CSI Parser
        ↓
Amplitude + Phase Feature Extraction
        ↓
Sliding Window Creation
        ↓
CNN + LSTM Deep Learning Model
        ↓
Real-Time Intrusion Prediction
\end{verbatim}

\section{Linux / Ubuntu Environment Setup}

\subsection{Step 1: Create Project Directory}

\begin{lstlisting}[language=bash]
mkdir -p ~/esp_projects
cd ~/esp_projects
\end{lstlisting}

\subsection{Step 2: Create Python Virtual Environment}

\begin{lstlisting}[language=bash]
python3.10 -m venv ~/esp_projects/csi_realtime_ml/csi_env
\end{lstlisting}

\subsection{Step 3: Activate the Environment}

\begin{lstlisting}[language=bash]
source ~/esp_projects/csi_realtime_ml/csi_env/bin/activate
\end{lstlisting}

After activation, the terminal should show:

\begin{verbatim}
(csi_env)
\end{verbatim}

\subsection{Step 4: Upgrade pip}

\begin{lstlisting}[language=bash]
python -m pip install --upgrade pip
\end{lstlisting}

\subsection{Step 5: Install Required Packages}

\subsubsection{Core Scientific Packages}

\begin{lstlisting}[language=bash]
pip install numpy pandas scipy matplotlib seaborn
\end{lstlisting}

\subsubsection{Machine Learning Packages}

\begin{lstlisting}[language=bash]
pip install scikit-learn tensorflow keras
\end{lstlisting}

\subsubsection{Serial Communication Packages}

\begin{lstlisting}[language=bash]
pip install pyserial
\end{lstlisting}

\subsubsection{Notebook and Kernel Packages}

\begin{lstlisting}[language=bash]
pip install notebook jupyterlab ipykernel
\end{lstlisting}

\subsubsection{Visualization Packages}

\begin{lstlisting}[language=bash]
pip install plotly
\end{lstlisting}

\subsubsection{Install All Packages Together}

\begin{lstlisting}[language=bash]
pip install \
numpy \
pandas \
scipy \
matplotlib \
seaborn \
scikit-learn \
tensorflow \
keras \
pyserial \
notebook \
jupyterlab \
ipykernel \
plotly
\end{lstlisting}

\section{Windows Environment Setup}

\subsection{Step 1: Create Environment}

\begin{lstlisting}[language=bash]
python -m venv csi_env
\end{lstlisting}

\subsection{Step 2: Activate Environment}

\begin{lstlisting}[language=bash]
.\csi_env\Scripts\activate
\end{lstlisting}

\subsection{Step 3: Install Packages}

\begin{lstlisting}[language=bash]
pip install numpy pandas scipy matplotlib seaborn
pip install scikit-learn tensorflow keras
pip install pyserial notebook jupyterlab ipykernel
pip install plotly
\end{lstlisting}

\subsection{Step 4: Register Jupyter Kernel}

\begin{lstlisting}[language=bash]
python -m ipykernel install \
--user \
--name csi_env \
--display-name "Python (csi_env)"
\end{lstlisting}

\section{Registering the Jupyter Kernel}

\subsection{Linux Kernel Registration}

\begin{lstlisting}[language=bash]
source ~/esp_projects/csi_realtime_ml/csi_env/bin/activate
\end{lstlisting}

\begin{lstlisting}[language=bash]
python -m ipykernel install \
--user \
--name csi_env \
--display-name "Python (csi_env)"
\end{lstlisting}

\subsection{Verify Kernel Installation}

\begin{lstlisting}[language=bash]
jupyter kernelspec list
\end{lstlisting}

Expected output:

\begin{verbatim}
csi_env
\end{verbatim}

\section{Opening the Project in VS Code}

\subsection{Open VS Code}

\begin{lstlisting}[language=bash]
cd ~/esp_projects/csi_intrusion_pipeline
code .
\end{lstlisting}

\subsection{Open Notebook}

Open:

\begin{verbatim}
Notebook/Train and test.ipynb
\end{verbatim}

\subsection{Select Kernel}

Inside the notebook:

\begin{verbatim}
Top Right -> Select Kernel
\end{verbatim}

Choose:

\begin{verbatim}
Python (csi_env)
\end{verbatim}

\subsection{If Kernel Does Not Appear}

Reload VS Code:

\begin{verbatim}
Ctrl + Shift + P
Developer: Reload Window
\end{verbatim}

Then select:

\begin{verbatim}
Python (csi_env)
\end{verbatim}

\section{Testing the Environment}

Run the following notebook cell:

\begin{lstlisting}[language=Python]
import numpy as np
import pandas as pd
import matplotlib.pyplot as plt
import sklearn
import serial
import tensorflow as tf

print("NumPy:", np.__version__)
print("TensorFlow:", tf.__version__)
print("GPU:", tf.config.list_physical_devices("GPU"))
\end{lstlisting}

Expected output:

\begin{verbatim}
GPU: [PhysicalDevice(name='/physical_device:GPU:0',
device_type='GPU')]
\end{verbatim}

\section{TensorFlow GPU Verification}

\subsection{Check NVIDIA GPU}

\begin{lstlisting}[language=bash]
nvidia-smi
\end{lstlisting}

\subsection{Check TensorFlow GPU Access}

\begin{lstlisting}[language=bash]
python -c "import tensorflow as tf; \
print(tf.config.list_physical_devices('GPU'))"
\end{lstlisting}

\section{Package Summary}

\begin{longtable}{|p{5cm}|p{8cm}|}
\hline
\textbf{Package} & \textbf{Purpose} \\
\hline
numpy & Numerical computation and array processing \\
\hline
pandas & CSV and dataframe handling \\
\hline
scipy & Signal processing utilities \\
\hline
matplotlib & Visualization and live plotting \\
\hline
seaborn & Advanced statistical visualization \\
\hline
scikit-learn & Machine learning preprocessing and evaluation \\
\hline
tensorflow & Deep learning framework \\
\hline
keras & CNN+LSTM model API \\
\hline
pyserial & ESP32 serial communication \\
\hline
notebook & Jupyter notebook support \\
\hline
jupyterlab & Interactive notebook interface \\
\hline
ipykernel & Jupyter kernel registration \\
\hline
plotly & Interactive visualization support \\
\hline
\end{longtable}

\section{Running the Notebook}

\subsection{Launch Jupyter Notebook}

\begin{lstlisting}[language=bash]
jupyter notebook
\end{lstlisting}

or

\begin{lstlisting}[language=bash]
jupyter lab
\end{lstlisting}

\section{Real-Time CSI Dashboard Execution}

\subsection{Activate Environment}

\begin{lstlisting}[language=bash]
source ~/esp_projects/csi_realtime_ml/csi_env/bin/activate
\end{lstlisting}

\subsection{Run Live Dashboard}

\begin{lstlisting}[language=bash]
cd ~/esp_projects/csi_intrusion_pipeline/realtime

python live_dashboard.py
\end{lstlisting}

\section{Common Problems and Fixes}

\subsection{Problem: No module named numpy}

This means the notebook is not using the correct kernel.

Fix:

\begin{verbatim}
Top Right -> Select Python (csi_env)
\end{verbatim}

\subsection{Problem: Python (csi\_env) does not appear}

Run:

\begin{lstlisting}[language=bash]
source ~/esp_projects/csi_realtime_ml/csi_env/bin/activate

python -m ipykernel install \
--user \
--name csi_env \
--display-name "Python (csi_env)"
\end{lstlisting}

Then reload VS Code.

\subsection{Problem: No module named serial}

Install pyserial:

\begin{lstlisting}[language=bash]
pip install pyserial
\end{lstlisting}

\subsection{Problem: TensorFlow GPU not detected}

Check GPU:

\begin{lstlisting}[language=bash]
nvidia-smi
\end{lstlisting}

Then check TensorFlow:

\begin{lstlisting}[language=bash]
python -c "import tensorflow as tf; \
print(tf.config.list_physical_devices('GPU'))"
\end{lstlisting}

\section{Final Recommended Kernel}

Always use:

\begin{verbatim}
Python (csi_env)
\end{verbatim}

for:

\begin{itemize}
    \item CSI parsing
    \item Feature extraction
    \item CNN+LSTM model training
    \item Model testing
    \item Real-time CSI inference dashboard
\end{itemize}

\end{document}
